\subsection{JPU 硬件 JPEG 解码}\label{sec:jpu}

机器人用的 USB 相机有两种交付像素的方式。一种是未压缩的 YUYV，整帧逐像素送出，占满带宽；
另一种是 MJPEG，即把视频的每一帧各自压成一张独立的 JPEG 图片再传出，体积小得多，因而是这类
相机在较高分辨率与帧率下的原生工作模式。代价是接收端拿到的每一帧都还是一张压缩的 JPEG，必须
先解码还原成像素才能进入后续处理。若由 CPU 来做这件事，一帧的 JPEG 解码在本平台上约需二十五
毫秒，几乎与一次推理相当，整条流水线会被这一步拖垮。RK3588 为此集成了一个专用的硬件 JPEG
解码单元，即 VDPU720（设备树节点 \texttt{jpegd@fdb90000}，本文称其为 JPU），它在硬件中完成熵
解码与反变换，直接把一张 JPEG 还原为 NV12 排布的 YUV 帧并写入物理内存，使这一步从计算路径上
让出。

在 Linux 上，应用并不直接操作这块硬件，而是经由瑞芯微的多媒体处理平台 MPP（Media Process
Platform）库。MPP 把"解码一帧"抽象成一次会话，应用通过 \texttt{librockchip\_mpp} 提交压缩
数据、取回解码帧，库再经字符设备 \texttt{/dev/mpp\_service} 把寄存器配置投向内核。\textbf{本队在 Starry
上重建了 \texttt{/dev/mpp\_service} 这一节点，实现了 \texttt{librockchip\_mpp} 的 JPEG 解码路径
所用的那一小部分 ioctl 约定，使未经改动的 MPP 程序，包括其标准解码测试 \texttt{mpi\_dec\_test}
与 gstreamer、ffmpeg 的瑞芯微后端，都能照常驱动这块解码器。}平台无关的驱动核心做成一个
\texttt{\#![no\_std]} 的 crate，负责把寄存器文件、量化表与霍夫曼表的字节排布按厂商 \texttt{hal\_jpegd\_rkv}
的方式拼出来；设备节点这一层只做用户态数据的拷入拷出、把 dma-buf 文件描述符翻译成物理地址，以及
真正驱动一次硬件运行。

\repfig{jpu-architecture.png}{相机 MJPEG 帧经 \texttt{librockchip\_mpp} 与 \texttt{/dev/mpp\_service}
进入 JPU 硬件解码，输出 NV12 帧落在与 RGA、NPU 共享的同一段 dma-buf 内存上}{fig:jpu}

\paragraph{对齐 MPP 的会话协议}
MPP 与内核之间并非一次函数调用，而是一串按字节对齐、经 ioctl 传入的请求记录。每条记录是一个
24 字节的结构，带有命令号、标志位与一个指向各自负载的用户态指针，多条记录以 \texttt{MULTI\_MSG}
标志串成一批，到 \texttt{LAST\_MSG} 为止。会话先以 \texttt{PROBE\_HW\_SUPPORT} 询问本节点支持
哪类编解码，节点据实回报只支持 JPEG 解码；再以 \texttt{QUERY\_HW\_ID} 读硬件版本，并以
\texttt{INIT\_CLIENT\_TYPE} 把会话绑定到 JPEG 解码这一客户端类型，绑定到其它类型会被拒绝。
随后 \texttt{SET\_REG\_WRITE} 送来要写入硬件的整组寄存器，\texttt{SET\_REG\_READ} 声明完成后要读
回的寄存器窗口，最后 \texttt{POLL\_HW\_FINISH} 触发一次提交并阻塞到这一帧解码完成。节点刻意不去
提供查询命令支持表的接口，从而让 MPP 走它对旧内核的兼容回退路径，恰好对上本节点实现的这一组
命令，使既有的 MPP 二进制无需任何改动。

\paragraph{把文件描述符翻译成解码器看得见的地址}
MPP 在寄存器组里并不直接写物理地址，而是把缓冲的 dma-buf 文件描述符放进几个约定的地址寄存器槽，
另以 \texttt{SET\_REG\_ADDR\_OFFSET} 附带每个槽的字节偏移。提交前节点逐一扫描这些地址槽，按厂商
\texttt{trans\_tbl\_jpgdec} 给定的位置，把描述符经共享的 \texttt{/dev/dma\_heap} 解析为其连续缓冲的
物理基址，再加上声明的偏移写回寄存器。量化表、霍夫曼最小码表与霍夫曼值表共处一段表缓冲，分别落在
偏移 384 与 704 字节处，输入码流与输出帧各占一个槽，于是 MPP 用文件描述符表达的一组缓冲被还原成
解码器能直接取用的物理地址。

\paragraph{让引擎在不开启地址翻译的情况下取到数据}
这块解码器的前面本来挂着一个属于它自己的小型地址翻译部件，即 IOMMU，作用是把设备发出的虚拟地址
翻译成内存中的物理地址。开启它需要为设备维护一套页表，而本流水线里所有缓冲都来自
\texttt{/dev/dma\_heap}，在四吉字节以下连续分配，物理地址本身就是连续的，无需翻译。因此设备初始化
时把这个 IOMMU 直接置为直通：它的寄存器在同一寄存器页偏移 \texttt{0x480} 处，先写入强制复位命令清掉
前一个系统可能遗留的页表，再把页表基址寄存器 \texttt{RK\_MMU\_DTE\_ADDR} 清零、写入停用分页命令
\texttt{RK\_MMU\_CMD\_DISABLE\_PAGING}，最后把它的中断全部屏蔽。此后引擎按物理地址直接读写内存，
\texttt{dma\_addr} 与物理地址、设备可见地址三者相等。

\paragraph{先配置寄存器，再启动引擎}
解码器的启动位与状态位同处控制寄存器 \texttt{SWREG1}。运行时把整组寄存器写入硬件时，刻意跳过只读的
ID 寄存器，并把带启动位的 \texttt{SWREG1} 留到最后一个写，从而保证几何、格式、表长与五个地址槽都
就位之后引擎才被拉起，绝不会在输入尚未配齐时就开始取数。完成与否由轮询 \texttt{SWREG1} 判断，
其中分别有解码就绪位与总线错误、码流错误、超时、码流为空几个错误位；读回整组寄存器供诊断之后，无论
成功、出错还是超时，都以写一清零的方式清掉这些状态位，使陈旧的完成或错误标志不会被下一帧误读。完成
信号这里以轮询观察，无需真正布线中断。

\paragraph{三个加速器共享同一段内存}
JPU 在本流水线里不孤立工作，它的输出正是下一级 RGA 色彩转换与缩放的输入，而 RGA 的结果又直接
写入 NPU 的输入缓冲。这要求三者读写同一段物理内存而非各自拷贝。承载解码输出的 NV12 帧由
\texttt{/dev/dma\_heap} 在四吉字节以下连续分配，以 dma-buf 文件描述符在 JPU、RGA、NPU 之间传递，
JPU 解出的帧落在这段缓冲上即可由 RGA 按描述符直接读取，整条 \texttt{MJPEG\,→\,JPU\,→\,RGA\,→\,NPU}
通路因此全程零拷贝。解码器是 32 位寻址，节点在解析文件描述符时拒绝四吉字节以上的缓冲而非默默截断
地址，恰与 dma-heap 的分配区间一致。

驱动核心的寄存器组、量化表与霍夫曼表拼装在主机上有完整的单元测试，包括对照厂商布局逐字节比对表
缓冲、断言启动位写在最后、以及量化表的反之字形重排。在开发板上，标准的 MPP 解码测试
\texttt{mpi\_dec\_test} 解出的结果与 Linux 逐字节一致，校验和为 \texttt{2712426383}，单帧解码的中位
延迟约一毫秒。需要说明的是，JPU 的意义并不在于比 YUYV 色彩转换路径更快，二者并不可比；它的意义在于
为相机的原生 MJPEG 模式提供一条硬件解码通路，使这一步不必回落到约二十五毫秒的 CPU 软解。接入之后，
\texttt{MJPEG\,→\,JPU\,→\,RGA\,→\,NPU} 全程零拷贝，端到端中位延迟为 9.06\,ms，约合每秒二十九帧，此时
的瓶颈已是相机供帧速率，与 YUYV 路径的 9.65\,ms 处于同一水平。\textbf{该工作以合并请求 \href{https://github.com/rcore-os/tgoskits/pull/1456}{\#1456} 提交至上游。}
